\title{Surface--like pseudolattices, del Pezzo surfaces, and genus one
fibrations}
\author{\textit{Based on a lecture series by }Franco Rota.}
\date{December 2024}
\maketitle

\subsection{About} This course on pseudolattices was a part of a Winter
School for students and early career researchers, organised by the UK Algebraic
Geometry Network. The lectures were delivered in-person in the University of
Warwick, and were transcribed by Parth Shimpi. These notes have undergone
several amendments and are not a verbatim recall of the lectures, therefore
discretion is advised when using this material. They are available online at
\url{https://pas201.user.srcf.net/documents/2023-ukag-pseudolattices.pdf}.
All errors and corrections should be communicated to by email to
\url{parth.shimpi@glasgow.ac.uk}.

Think about homological mirror symmetry on a linear algebraic level.

We recall some of the bilinear--algebraic constructions associated to a smooth
projective \(\mathbb{C}\)--surface \(S\). A \emph{Cartier divisor} \(D\) on
\(S\). To such \(D\) we can associate a line bundle \(\mathscr{O}_S(D)\), and
this gives a one--to--one correspondence between the set of divisors on \(S\) up
to rational equivalence and the set of line bundles on \(S\) up to isomorphism.
Either of these is a group called the \emph{Picard group}.

\subsection{The bilinear pairing on the Picard group} Given \(C,C'\) distinct
irreducible curves in \(S\) and a point \(x\in C\cap C'\), one can associate an
integer \(\text{mult}_x(C\cap C')\) called the \emph{intersection multiplicity}
of \(C,C'\) at \(x\). If \(C,C'\) are locally defined by polynomials \(f,f'\)
respectively, then the intersection multiplicity is the dimension of
\(\mathscr{O}_x/(f,f')\). Then the intersection pairing is defined as
\[C\cdot C' = \sum_{x\in C\cap C'} \text{mult}_x(C\cap C')\]
and we extend it linearly to sums of curves.

Fact:\(\mathscr{O}_{C\cap C'}=\mathscr{O}_S/\mathscr{O}_S(-C) +
\mathscr{O}_S(-C')\), so \(C\cdot C'=\hh^0(S,\mathscr{O}_{C\cap C'})\).

\begin{theorem}
For \(L,M\in \Pic(S)\),
\[
    (L\cdot M) = \chi(\OO_S) - \chi(L\dual) - \chi(M\dual) +\chi(L\dual\otimes
    M\dual)
\]
is a symmetric bilinear form that extends \((C\cdot C') = (\OO_S(C)\cdot
\OO_S(C'))\).
\end{theorem}

Topological interpretation: since \(S\) is a complex variety, there is an
exponential exact sequence
\[0\to \mathbb{Z}\to \OO_S \xrightarrow{e^{2\pi i (-)}} \OO_S^\times \to 0.\]
Then \(\Pic(S)\) is isomorphic to \(\HH^1(S,\OO_S^\times)\), so taking the long
exact sequence in cohomology gives
\[...\to \HH^0(S,\OO_S)\to \HH^0(S,\OO_S^\times) \to \HH^1(S,\mathbb{Z})\to
\HH^1(S,\OO_S)\to \Pic(S)\xrightarrow{c_1} \HH^2(S,\mathbb{Z})\to
\HH^2(S,\OO_S)\to ...\]
Now the map \(\HH^0(S,\OO_S)\to\HH^0(S,\OO_S^\times)\) is simply the exponential
map \(\mathbb{C}\to \mathbb{C}^\times\) and is in particular surjective. Thus
the Picard group sits in an exact sequence
\[
0\to \frac{\HH^1(S,\OO_S)}{\HH^1(S,\mathbb{Z})} \to \Pic(S) \xrightarrow{c_1}
\text{NS}(S)\to 0
\]
where \(\text{NS}(S)\subset \HH^2(S,\mathbb{Z})\) is the image of
\(c_1:\Pic(S)\to \HH^2(S,\mathbb{Z})\). We say \(\text{NS}(S)\) is the
\emph{N\'eron--Severi group} of \(S\), and the map \(c_1\) is the first
\emph{Chern class} map. Given an irreducible curve \(C\subset S\), there is a
restriction map \(\HH^2(S,\mathbb{Z})\to \HH^2(C,\mathbb{Z})=\mathbb{Z}\) which
by Poincar\'e duality gives \(c_1(C)\in \HH^2(S,\mathbb{Z})\).

Then by construction, the cup product on cohomology gives the intersection product \(C\cdot
C'=c_1(C)\cap c_1(C')\).

\subsection{Chern characters} If \(X\) is a smooth projective variety, then each
line bundle \(L\in \Pic (X)\) can be assigned a \emph{Chern polynomial}
\(\cc(L)_t = 1+\cc_1(L)t \in \HH^\bullet(X,\mathbb{Z})[t]\). We say the
\emph{Chern character} is \[\ch(L)=e^{\cc_1(L)} =
1+\cc_1(L)+\frac{1}{2}\cc_1(L)^2+...\]
If \(E=\bigoplus_{i=1}^rL_i\) is a direct sum of line bundles (or more generally
if \(E\) is a vector bundle filtered by line bundles \(L_1,...,L_r\)), then the
Chern polynomial and character are defined as
\[
\cc(E)_t = \prod_{i=1}^r\cc(L_i)_t = 1 + \cc_1(E)_t + \cc_2(E)t^2 + ... \\
\ch(E) = \sum_{i=1}^r\ch(L_i) = \sum_{k=0}^{\dim X} \ch_k(E).
\]
More generally, the splitting principle lets us assume \(\cc(E)_t\) always has
the form \(\prod(1+a_it)\) for cohomology classes \(a_i\in \HH^2(X,\mathbb{Z})\)
called the Chern roots of \(E\). These can be used to define the \emph{Todd
class}
\[
\todd(E)= \prod_{i=1}^r \frac{a_i}{1-e^{-a_i}}.
\]

\begin{exercise}[Properties of Chern classes]
\(\ch(E\oplus F)=\ch(E)+\ch(F)\), \(\ch(E\otimes F)=\ch(E)\cdot \ch(F)\),
\(\ch(E\dual)=\sum_{i=1}^{\dim X}(-1)^i\ch_i(E)\).
\end{exercise}

\begin{exercise}[Chern and Todd classes on surfaces]
If \(E\) is a vector bundle on a surface \(S\), then
\[
\ch(E)=1+\cc_1(E)+\frac{\cc_1^2(E)-2\cc_2(E)}{2}, \qquad
\todd(E) = 1+\frac{1}{2} \cc_1(E) + \frac{1}{12} (\cc_1^2(E)+\cc_2(E)).
\]
\end{exercise}

\begin{theorem}[Hirzebruch--Riemann--Roch]
If \(E\) is a vector bundle on a smooth projective variety \(X\) of dimension
\(n\), then
\[\chi(E) = \sum_{i=0}^n (-1)^i \dim \HH^i(X,E) = \deg[\ch(E)\cdot \todd(T_X)]_n
= \int_X \ch(E)\todd(T_X)\].
\end{theorem}

\begin{example}
If \(X=C\) is a smooth curve and \(L=\OO_C(\sum n_i p_i)\) is a line bundle,
then \(\todd_C=\todd(T_C)=\todd(\omega_C\dual) = 1+\frac{1}{2}\cc_1(\omega\dual)
= 1+\frac{2-2g}{2}[p]\) while \(\ch(L) =1+(\sum n_i)[p]\). Then HRR gives us
\(\chi(L) = \deg(L)+1-g\).
\end{example}

\begin{example}
If \(X=C\) is a surface, then \(\todd_S=(1,-\frac{K_S}{2},
\frac{1}{12}(\cc_1^2+\cc_2))\) and \(\ch(L)=(1,L,\frac{L^2}{2})\). Thus HRR gives
us
\[\chi(L)=\underbrace{\frac{1}{12}(\cc_1^2+\cc_2)}_{\chi(\OO_S)} -
\frac{1}{2}(L\cdot K_S) + \frac{L^2}{2} = \chi(\OO_S)+\frac{1}{2}(L-K_S)\cdot L.\]
Noether's formula \(\chi(\OO_S)=\frac{1}{12}(\cc_1^2+\cc_2)\).
\end{example}

Extensions: If \(E,F\) are vector bundles, then have Euler pairing
\(\chi(E,F)=\sum_{i=0}^n (-1)^i \ext^i (E,F)\). Since \(\chi(E,F)=\chi(F\otimes
E\dual)\), HRR gives us \(\chi(E,F)=\int_X \ch(F)\ch(E\dual)\todd_X\).


Since \(X\) is smooth, each sheaf \(G\) has a finite resolution by vector
bundles \(0\to E_n\to E_{n-1}\to...\to E_0\to G\to 0\).  Define
\(\ch(G)=\sum_{i=0}^n(-1)^i \ch(E_i)\)

Define: Grothendieck group of an Abelian category, triangulated category,
derived category.

Upshot: \(\chi:\KK_0(\Coh X)\times \KK_0(\Coh X)\to \mathbb{Z}\) is a bilinear
form, not necessarily symmetric.

Define \(\ker\chi = \{G\;|\; \chi(G,H)=0\text{ for all }H\in\KK_0(\Coh X)\}\).
In the presence of a Serre functor, this is equal to  \(\{G\;|\;
\chi(H,G)=0\text{ for all }H\in\KK_0(\Coh X)\}\).

The intersection pairings on \(\Pic(X), \NerSev(X)\) are symmetric. On the other
hand, the Euler pairing on \(\KK_0(\Coh X)\) is not usually symmetric.

[Very rough introduction to derived category, cohomology functors,
quasi-isomorphisms and roofs, Ext1--Hom correspondence, replacing objects by
resolutions, Serre functors]

\begin{theorem}[Serre duality]
    Given a smooth projective variety \(X\) and complexes \(E,F\in \Db X\),
    there is a canonical isomorphism
    \[
        \Hom_{\Db X}(E,F) \cong \Hom_{\Db X}(F, E\otimes \omega_X[\dim X])\dual.
    \]
\end{theorem}

References: FM transforms in AG, Gelfand--Manin

\begin{theorem}
    Given a triangulated category \(\mathscr{T}\), a Serre functor is an
    autoequivalence \(\mathbb{S}:\mathscr{T}\to \mathscr{T}\) such that
    \(\Hom(A,B)\cong \Hom(B,\mathbb{S}(A))\dual\).
\end{theorem}

Suppose \(\mathscr{T}\) is triangulated with a Serre functor \(\mathbb{S}\).
Then \(\ker \chi = \{E\in \mathscr{T}\;|\; \chi(E,F)=0 \text{ for all
}F\in\mathscr{T}\}\) is equal to \(\{E\in\mathscr{T}\;|\; \chi(F,E)=0 \text{
for all }F\in\mathscr{T}\}\) since \(\chi(E,F)=\chi(\mathbb{S}^{-1}F, E)\). This
lets us define the numerical Grothendieck group \(\KK_0^\text{num}=\KK_0/\ker
\chi\).

\begin{example}
    For \(E\) an elliptic curve, \(\KK_0(E)\) is not finitely generated since it
    contains \(\Pic^0(E)\cong E\). On the other hand, \(\KK_0^\text{num}(E)\) is
    isomorphic to \(\mathbb{Z}^2\) with basis \([\OO_p]\), \([\OO_E]\). What's the
    Euler pairing?
\end{example}

\begin{definition}
    A \emph{pseudolattice} is a finitely generated free Abelian group \(G\)
    equipped with a non--degenerate bilinear pairing \(\chi:G\times G\to
    \mathbb{Z}\). We say \(G\) is \emph{unimodular} if the induced natural map
    \(\xi:G\to \Hom(G,\mathbb{Z})\) given by \(v\mapsto \chi(-,v)\)  is an
    isomorphism.

    We say pseudolattices \(G,G'\) are isometric if there is a group isomorphism
    \(\varphi:G\to G'\) that respects the pairings.
\end{definition}

\begin{definition}
    A \emph{Serre operator} on \(G\) is an isometry \(\mathbb{S}:G\to G\) such that
    \(\chi(u,v)=\chi(v,\mathbb{S}u)\) for all \(u,v\in G\).
\end{definition}

\begin{exercise}
    Defining \(\xi^T:G\to G\dual\)by \(v\mapsto \chi(v,-)\), show that a Serre
    operator exists if and only if \(\xi\) is an isometry and in this case
    \(\mathbb{S}=\xi^{-1}\circ \xi^T\).
\end{exercise}

\subsection{Exceptional bases and mutation} An element \(e\in G\) is
(numerically) exceptional if \(\chi(e,e)=1\). A (numerical) exceptional sequence
is a sequence \((e_1,...,e_n)\) of exceptional elements such that
\(\chi(e_i,e_j)=0\) whenever \(j\leq i\). An exceptional basis is an exceptional
sequence which spans \(G\).

\begin{lemma}
    If there is an exceptional sequence of length \(\textnormal{rank}(G)\) then the
    sequence is an exceptional basis and \(G\) is unimodular.
\end{lemma}

Mutations: given \(e\in G\) exceptional, we can define operators
\[\mathbb{L}_e(v)=v-\chi(e,v)e, \quad \mathbb{R}_e(v)=v-\chi(v,e)e\]

Given an exceptional sequence \(e_\bullet=(e_0,...,e_n)\) and \(0\leq i \leq
n\), define \(\mathbb{L}_i(e_\bullet)=(...,e_{i-1}, \mathbb{L}_{e_i}(e_{i+1}),
e_i, e_{i+2}, ...)\) and  \(\mathbb{R}_i(e_\bullet)=(...,e_{i-1}, e_{i+1},
\mathbb{R}_{e_{i+1}}(e_i), e_i, e_{i+2}, ...)\).

Modelling on \(\KK_0^\text{num}(X^2)=\HH^0(X)\oplus \NerSev(X)/\text{tors} \oplus
\HH^4(X)\), we make the following definition.

\begin{definition}[Kuznetsov, Vial]
A pseudolattice \((G,\chi)\) is \emph{surface--like} if there is a primitive
element \(p\in G\) such that \(\chi(p,p)=0\), \(\chi(p,-)=\chi(-,p)\), and
\(\chi|_{p\orth}\) is symmetric.
\end{definition}

We say such a \(p\) is pointlike. For \(u,v\in p\orth\), write
\(\chi(u,v)=-u\cdot v\). Write \(\text{rk}(v)=\chi(v,p)\), then \(p\orth =
\{v\in G\;|\; \text{rk}(v)=0\}\).

There is a chain complex \(\mathbb{Z}\xrightarrow{p} G \xrightarrow{\text{rk}}
\mathbb{Z}\) with middle cohomology \(p\orth/p\) which we call the
\emph{Neron--Severi group} \(\NerSev(G)\). This is a finitely generated free
Abelian group of rank \(\text{rk}G-2\).

If \(G\) has a Serre operator, then the filtration \(0\into \mathbb{Z}p\into
p\orth \into G\) has backward maps \(G\xrightarrow{\text{id}-\mathbb{S}}p\orth
\xrightarrow{\text{id}-\mathbb{S}} \mathbb{Z}p
\xrightarrow{\text{id}-\mathbb{S}} 0\). In particular
\((\text{id}-\mathbb{S})^3=0\) i.e.\ \(\mathbb{S}\) is unipotent.

\subsection{The canonical class} If \((G,\chi)\) is a surface--like
pseudolattice, then there is a unique object \(K\in
\text{NS}(G)\otimes\mathbb{Q}\) such that for any \(v_1,v_2\in G\),
\(\chi(v_1,v_2)-\chi(v_2,v_1)= -K \cdot \lambda(v_1,v_2)\)  where
\(\lambda(v_1,v_2)=\text{rk}(v_1)v_2 - \text{rk}(v_2)v_1)\).

In the surface setting, \(K=K_X\) is the class of a canonical divisor.

The motivation to consider numerical exceptional collections is to obstruct
the existence of exceptional collections in the derived category.

\begin{definition}
    If \(\mathscr{T}\) is a triangulated category, then \(E\in \mathscr{T}\) is
    exceptional if \(\Hom^\bullet(E,E)=\mathbb{C}\). An exceptional collection
    is a sequence \((E_1,...,E_n)\) of exceptional objects such that
    \(\Hom(E_i,E_j[k])=0\) for all \(i>j\), and \(k\in \mathbb{Z}\).


    We say it is \emph{full} if it generates \(\mathscr{T}\).
\end{definition}

It is immediate that the classes \([E_1],...,[E_n]\) form an exceptional basis
of \(\KK_0^\text{num}\).

The leading conjecture (by Orlov) is that a smooth projective surface over
\(\mathbb{C}\) admits a full exceptional collection if and only if it is
rational.

Other questions:
\begin{enumerate}
    \item Given an exceptional collection of maximal length
        (\(2+\text{rk}(\NerSev X)\)), when is it full?
    \item When do mutations act transitively on the set of all exceptional
        collections of maximal length?
\end{enumerate}

Answers: for del Pezzo surfaces (i.e.\ blow ups of \(\mathbb{P}^2\) in \(\leq
8\) points in general position), the answer to both the questions is `always'
(see Pirozhkov `20, Kuleshov--Orlov `94). Also for blow up in 9 points [Krah].

[Cheng--Hayden--Schroll] consider the Fukaya category of a specific Riemann
surface to show that the mutation action need not be transitive.

[Krah] found that the Blow up of \(\mathbb{P}^2\) in ten points does \emph{not}
have a full exceptional collection of maximal length, thus having at least two
orbits.

We focus on numerical versions of these.

\begin{theorem}[(Vial)]
    If \(X\) is is a smooth projective surface with \(\chi(\OO_X)=1\) i.e.\
    \([\OO_X]\) is numerically exceptional, then the following are
    equivalent.
    \begin{enumerate}
        \item \(X\) has a numerical exceptional sequence \(L_0,...,L_{n+1}\) of
            line bundles of maximal length.
        \item \(X\) has a numerical exceptional sequence \(E_0,...,E_{n+1}\) of
            numerical line bundles (i.e.\ \(\text{rk}(E_i)=1\)) of maximal
            length.
        \item \(K_X^2=10-\rho(X)\) and one of the following holds
            \begin{enumerate}
                \item \(\mathbb{P}^2\)--like: \(\NerSev(X)=\langle 1 \rangle\)
                    and \(\frac{1}{3}K_X\) is primitive.
                \item \(\mathbb{P}^1\times \mathbb{P}^1\)--like:
                    \(\NerSev(X)=U\) (hyperbolic plane) and \(\frac{1}{2}K_X\)
                    is primitive.
                \item Blow up of \(\mathbb{P}^2\) in \(n-1\) like:
                    \(\NerSev(X)=\langle 1 \rangle \oplus \langle -1 \rangle
                    ^{n-1}\) and \(K_X\) is primitive.
            \end{enumerate}
    \end{enumerate}
\end{theorem}

Main tool [Hille--Perling] is to attach a toric system-- given
\((E_0,...,E_{n+1})\), take \(D_i=\cc_1(E_i)-\cc_1(E_{i-1})\) to get a set
\(D_1,...,D_n\) which forms a trigonal unimodular basis for \(\NerSev(X)\). If
you assume in addition that the signature is \((1,n-1)\) and there is a
canonical class (characteristic element) then such lattices have been completely
classified.

Geometric consequences: if \(X\) is a smooth projective surface, \(\OO_X\) is
exceptional if and only if \(p_g=\hh^0(K_X)\) and \(q=\HH^1(\OO_X)\) both
vanish. Such surfaces have been classified by [Kotschik] according to Kodaira
dimension.

For \(k=-\infty\) these are the minimal surfaces, i.e.\ \(\mathbb{P}^2\),
\(\mathbb{P}^1\times \mathbb{P}^1\), and Hirzebruch surfaces \(\bbF_n\) for
\(n\geq 2\).

If \(k=0,1\) these are the Dolgachev surfaces \(X_9\) (Blow up of
\(\mathbb{P}^2\) in \(9\) points) the \(9\) points form the base locus of a
pencil of cubics thus giving an elliptic fibration \(X_9\to \mathbb{P}^1\). You
can also apply a log transformation to get \(X_9(p_1,...,p_\ell)\) where you add
multiplicity \(p_i\) to a smooth fiber \(F_i\). If \(k=0\) this gives the
Enriques surface \(X_9(2,2)\), all others have \(k=1\).

For \(k=2\) these are some fake projective planes (including Godeaux surfaces).

\begin{theorem}[(Vial)]
    \(X\) smooth projective surface with \(p_g=q=0\). If \(S\) is not minimal
    then there is a numerically exceptional collection of maximal length. If
    \(S\) is minimal then there is a numerically exceptional collection of
    maximal length except for
    \(X_9(2,2),X_9(2,3),X_9(2,4),X_9(3,3),X_9(2,2,2)\).
\end{theorem}

Enriques surfaces therefore do not admit exceptional collections. [Cho] gives an Enriques surface whose blowup has a full exceptional collection.

Kuznetsov recovers a similar classification result-- if \(G\) is a surface--like
pseudolattice of rank \(n+2\), we say it is geometric if \(\NerSev(G)\) has
signature \((1,n-1)\). We say \(G\) is minimal if there are no exceptional
objects of rank \(0\). Then Kuznetsov sets up a `numerical MMP': if \(e\in G\)
is an exceptional object of rank \(0\), then \(G_e=\orth e\) and have a
semi--orthogonal decomposition \(G=e\oplus G_e\) with pairing
\(\begin{pmatrix}1&\ast\\0&\chi_{G_e}\end{pmatrix}\) and
\(\NerSev(G)=\NerSev(G_e)\oplus \langle 1 \rangle\) is an orthogonal decomposition.
Then induction and toric systems game recovers the classification.

\section{The Fano/LG correspondence} Start with a Fano variety/orbifold \(X\).
The mirror is an LG model \(W:Y\to \mathbb{C}\) where \(Y\) is a non--compact
manifold and \(W\) is smooth. Mirrors to surfaces are a source of surface--like
pseudolattices-- \((\KK_0^\text{num}(FS(Y,W)),\chi)\) is surfacelike.

But theres more structure-- if there is a smooth effective \(D\in |-K_X|\)
(quasi del Pezzo: \(X\) a rational surface with smooth effective anticanonical
divisor, which is necessarily an elliptic curve by adjunction) then \(i:D\into
X\) gives linear map \(i^\ast:\KK_0^\text{num}X\to \KK_0^\text{num}D\)
induced by a functor \(i^\ast: \Db X \to \Db D\) which is spherical. Main
reference: [Harder--Thompson]

\begin{definition}[Addington, Anmo--Logvinenko]
    A spherical functor is an exact functor \(F:\mathcal{A}\to \mathcal{B}\) of
    triangulated categories with left and right adjoints \(L,R:B\to A\) such
    that the cotwist functor \(C\) defined as \(C\to \text{id}\to RF\) is an
    equivalence and \(R\cong CL\).
\end{definition}

If \(F\) is spherical then we can also prove that the twist \(T\) defined as
\(FR\to \text{id}\to T\) is also an equivalence.

\begin{example}
    An object \(E\in \mathscr{B}\) is \(n\)--spherical if \(\mathbb{S}(E)=E\)
    and \(\Hom(E,E[i])=\mathbb{C}\oplus \mathbb{C}[n]\).
    Then the map \(\Db(\ast)\to \mathscr{B}\) given by \(V\mapsto V\otimes E\)
    is a spherical functor.

    Why? Note that \(RF=\text{id}\oplus [-n]\) so the Cotwist is \([-n-1]\).

    Thus also the twist functor is an equivalence, it is given by the cone of
    \(\Hom(E,X)\otimes E \to X\).
\end{example}

If \(D\subset X\) is a smooth anticanonical divisor, then \(i_\ast\) is right
adjoint to \(i^\ast\),  and the cotwist is \(-\otimes \OO(-D)=-\otimes\omega_X\).

[Harder--Thompson] if \(f:G\to H\) is a Hom of pseudolattices, \(f\) is
spherical if there is a right adjoint \(r:H\to G\) and \(T=\text{id}-fr,
C=\text{id}-rf\) are invertible.

A spherical functor \(f:G\to H\) is relative \((-1)^n\)--CY if \(C=(-1)^n\mathbb{S}\).

If \(X\) is a quasi--del Pezzo and \(D\) is an elliptic curve,
\(\mathbb{E}=\KK_0^\text{num}(D)\) in the skyscraper--structuresheaf basis
\(\{a,b\}\) is given by the lattice \((\begin{smallmatrix} 0 & -1 \\ 1 &
0\end{smallmatrix})\); this is isomorphic to the lattice \(\HH^1(D,\mathbb{Z})\)
with intersection pairing.

We say \(f:G\to \mathbb{E}\) is quasi--del Pezzo if \(f\) is a relative 0-CY
spherical with adjoint \(r\), \(G\) is surfacelike with
pointlike \(p=r(a)\), and \(G\) admits an exceptional basis \(e_\bullet\) with
primitive images in \(\mathbb{E}\), and \(\NerSev(G)\) has sign \((1,n-1)\).

Can then show that \(a\) is fixed by twist and \(T\) acts on \(\mathbb{E}\) by \((\begin{smallmatrix} 1&-K_G^2\\0&1 \end{smallmatrix})\), further \(r(b)\in p\orth\) and \(r(b)=-K_G\) in \(\NerSev(G)\)

Notation: given \(\mathbb{E}=\langle a,b \rangle\) and \(v_i\in \mathbb{E}\),
can define a pseudolattice \(Z(\{v_i\})\) with formal basis \(z_i\) and
\(\chi(z_i,z_j)=\langle v_i,v_j \rangle_\mathbb{E}\).

Classification: up to mutations of \(e_\bullet\) and automorphisms
of \(\mathbb{E}\) of the form \((\begin{smallmatrix}1&K\\0&1\end{smallmatrix})\)
any quasiDelpezzo is of the form
\[
    (A)\quad Z(b,3a+b, 6a+b, a,...,a)\to \mathbb{E}\qquad
    (B)\quad Z(b,2a+b,2a+b,4a+b)\to \mathbb{E}
\]

Exercise:what are the surfaces that realise these cases?

Exercise: show that there is an isomorphism \(Z(b,3a+b,6a+b,a,a)\cong
Z(b,2a+b,2a+b,4a+b,a)\)

What does this story look like on the mirror? Genus 1 Lefschetz fibrations
\(W:Y\to \mathbb{C}\) with generic fiber given by a torus and special fibers given
by pinched tori. Choose paths from critical points to a point \(\infty\in
\mathbb{C}\), lift (parallel transport) to get vanishing cycles \(Z_i\in
F_\infty\). These define cohomology classes \(V_i\in
\HH^1(F_\infty,\mathbb{Z})\cong \mathbb{E}\). Then the vanishing cycles generate
\(\text{FS}(Y,W)\) and so \(\KK_0^\text{num}(FS(Y,W))\) is isomorphic to
\(Z(\text{vanishing thimbles }v_i)\cong \HH_2(Y,F\;;\mathbb{Z})\). By the
Harder--Thompson classification, get `numerical mirror symmetry
statements:\(\KK_0^\text{num}(X)\to \KK_0^\text{num}(D)\) is isomorphic to
\(\KK_0^\text{num}(FS(Y,W))\to \HH^1(F,\mathbb{Z})\) and further the isomorphism
is given by sequences of mutations. Trick is to take an \(e_\bullet\) that is
minimal in the sense of \(\sum |\text{rk}e_i|^2\).

Question: To prove that \(\KK_0^\text{num}(FS(Y,W))\cong Z(\text{vanishing
thimbles})\) we use \(g=1\) fibers. Why doesn't this work for higher genus? What
if you take fibrations that aren't Lefschetz? What if \(X\) is singular (log dP)
and \(|-K_X|=\emptyset\) but \(|-2K_X|\) has sections? McKay correspondence says
you should look at an orbifold inside which \(\Db(\tilde X)\) lives.





